\chapter{Литературный priority audit mapping-cone spectral Hessian}

По состоянию на 13 августа 2026 года выполнен targeted literature search по
четырём связкам ключевых слов: block log-determinant Hessian и Schur
complement; determinant line и mapping-cone torsion; fidelity susceptibility
и Wilson loops; graph incidence и pullback log-det geometry. Цель аудита ---
не объявить приоритет, а отделить известные общие механизмы от конкретной
специализации текущей программы.

\section{Что в литературе уже существует}

\begin{enumerate}
 \item \textbf{Mapping-cone determinant lines.}
 Kaad и Nest, arXiv:1403.7937, исследуют torsion isomorphisms determinant
 lines mapping-cone triangles Fredholm complexes и их canonical holomorphic
 sections. Следовательно, использование mapping cone как носителя
 determinant geometry не является новым само по себе.

 \item \textbf{Log-det Hessian как pullback metric на edge space.}
 В arXiv:2603.24913 Hessian функции \(-\log\det X\) явно pullback-ится на
 пространство matrix/graph edge weights:
 \begin{equation}
  g(U,V)=\Tr\bigl(X^{-1}L(U)X^{-1}L(V)\bigr).
 \end{equation}
 Это очень близкий общий шаблон: incidence deformation входит в block
 operator, а inverse-square response возникает из log-det Hessian.

 \item \textbf{Inverse-gap-square quantum geometry.}
 Zanardi--Giorda--Cozzini (\path{quant-ph/0701061}) выводят
 natural metric/fidelity susceptibility на parameter space Hamiltonians.
 К тому же выводу приходят Campos Venuti--Zanardi
 (\path{arXiv:0705.2211}). Spectral representation
 такой metric содержит squared inverse gaps. Поэтому интерпретация Hessian
 как response geometry также имеет прямых предшественников.

 \item \textbf{Schur complement и Gaussian determinant.}
 Block Gaussian models, conditional covariances и effective actions
 стандартно редуцируются Schur complement и log determinant. Этот
 algebraic механизм широко известен и не является самостоятельной заявкой
 на новизну.

 \item \textbf{Wilson loops рядом с quantum geometry.}
 В band-topology и quantum-criticality литературе Wilson loops и quantum
 geometric tensor используются в одной задаче, например в arXiv:2312.01086.
 Однако там Wilson loop служит topological diagnostic; он не выступает
 mapping-cone incidence внутри zeta-logdet Hessian с коэффициентом
 \(2\zeta(4)/\pi^4\).
\end{enumerate}

\section{Что при поиске не обнаружено}

Не найдена работа, где одновременно присутствуют все следующие элементы:
\begin{enumerate}
 \item periodic operator со spectrum \((\pi n)^2\), \(n\ne0\);
 \item rank-one или mapping-cone coupling, построенный из
 \(b_\theta=(1-U_\theta)u\);
 \item real-Gaussian factor \(1/2\) и Schur block
 \(\mathbb R\oplus\mathbb R^2\);
 \item Hessian coefficient
 \begin{equation}
  -\frac12\|b_\theta\|^2
  \Tr{}'\widehat\Delta^{-2}
  =-\frac{1-\cos\theta}{45};
 \end{equation}
 \item интерпретация этого coefficient как missing Wilson periodic
 primitive при одновременном запрете использовать full determinant из-за
 higher powers.
\end{enumerate}

Поиск по точным сигнатурам \(1/45\), \(\zeta(4)\),
\(1-\cos\theta\), ``mapping cone'', ``Wilson'', ``fidelity
susceptibility'' и ``log determinant Hessian'' не дал совпадающей
конструкции. Это отсутствие найденного совпадения, а не математическое
доказательство отсутствия публикации.

\begin{proposition}[Осторожный priority verdict]
Общие ингредиенты mapping-cone spectral Hessian известны и имеют сильные
литературные предшественники. Потенциально новой выглядит не сама техника, а
её конкретная specialization, в которой primitive Wilson incidence и
periodic zeta spectrum совместно дают coefficient \(1/45\). До независимого
поиска в MathSciNet, zbMATH, INSPIRE, Crossref и консультации специалиста
нельзя утверждать научный приоритет.
\end{proposition}

\begin{protocol}[Правило будущего заявления]
В публичном статусе допустима формулировка:
``найдена, по-видимому, ранее не отмеченная specialization известных
log-det-Hessian и mapping-cone механизмов''. Формулировки ``впервые
доказано'' и ``никто ранее не замечал'' запрещены до внешнего priority audit
и рецензирования.
\end{protocol}
